\documentclass[11pt,a4paper]{article}

% 基本包
\usepackage[utf8]{inputenc}
\usepackage[hidelinks]{hyperref}
\usepackage{enumitem}
\usepackage{xcolor}

\usepackage{luatexja-fontspec}

% \setmainfont{Arial}                    % 设置英文无衬线字体
% \setsansfont{Arial}                    % 如果想为无衬线设置单独字体，可以用这行
\setmonofont{Courier New}              % 英文等宽字体
% \setCJKmainfont{Noto Sans CJK JP}      % 设置中文主字体为无衬线字体
% \setCJKsansfont{Noto Sans CJK JP}      % 中文无衬线字体
% \setCJKmonofont{Noto Sans Mono CJK JP} % 中文等宽字体
\setmainjfont{Noto Sans CJK JP}        % 主体日文/中文
\setsansjfont{Noto Sans CJK JP}        % 无衬线
\setmonojfont{Noto Sans Mono CJK JP}   % 等宽


\definecolor{darkblue}{RGB}{0,0,139}   % 深蓝
\definecolor{burgundy}{RGB}{128,0,32} % 暗红


% 高亮宏（你在原文中用到的）
% \newcommand{\highlight}[1]{\textcolor{blue}{\textbf{#1}}}
\newcommand{\highlight}[1]{\textcolor{burgundy}{\textbf{#1}}}
\newcommand{\strongit}[1]{\textbf{\textit{#1}}}

\title{}
\date{} % 去掉日期
\author{} % 不显示作者

\begin{document}

\begin{center}
  \Large \textbf{Overview of Five Representative Publications}
\end{center}
\vspace{0.5em}

\noindent \textbf{Applicant:} Deyun Lyu \\
\textbf{Affiliation and Position:} Research Fellow, National Institute of Informatics (NII), Tokyo, Japan\\
\textbf{Position Applied For:} Assistant Professor (Informatics and Data Science, Graduate School of Advanced Science and Engineering)\\
\textbf{Field of Specialization:} Intelligent Science (Cyber-Physical Systems, Artificial Intelligence, Software Engineering – with a focus on Quality Assurance)\\
\textbf{Summary:} The following five publications have been selected as my \textbf{representative works}, as they collectively establish a coherent research program on the \textbf{quality assurance of AI-enabled cyber-physical systems (AI-CPSs)}. They address the full lifecycle of quality assurance, covering \textit{benchmarking, testing, fault localisation, and repair}, and have been published in top-tier venues (ICSE, TSE, TOSEM, GECCO, JSS). These works not only advance academic foundations but also contribute to practice through collaborations with Toyota, GAIO, and the Japan Automobile Research Institute, where one representative deployment target is end-to-end autonomous driving systems powered by AI.

\section*{Five Representative Publications}

My doctoral research focuses on the \textbf{quality assurance of AI-enabled cyber-physical systems} (AI-CPSs), a critical challenge for ensuring the reliability and trustworthiness of AI controllers embedded in physical systems such as autonomous vehicles and robots. These five publications collectively span the lifecycle of quality assurance of such systems, including benchmarking [AI-CPS-P1], testing [AI-CPS-P2], fault localisation [AI-CPS-P3, AI-CPS-P4], and repair [AI-CPS-P5]. Prior to this, existing research often missed the forest for the trees (木を見て、森を見ず): (i) no systematic benchmarking or evaluation of AI-CPSs existed; (ii) testing focused only on system-level correctness while ignoring DNN behaviour; and (iii) fault localisation and repair targeted simple DNN classifiers, failing to capture system-level complexity. In contrast, my research \textbf{sees both the trees and the forest} (木も見て、森も見る), establishing a \textbf{system–component co-analysis} methodology that advances quality assurance and paves the way for trustworthy AI-CPSs. Beyond academic contributions, these techniques are being advanced into practice through the collaboration with Toyota and the Japan Automobile Research Institute, toward deployment in their end-to-end autonomous driving systems.

\begin{itemize}[leftmargin=*]
    \item \textbf{[AI-CPS-P1]} Jiayang Song, \textbf{Deyun Lyu}*, Zhenya Zhang, Zhijie Wang, Tianyi Zhang and Lei Ma. ``\href{https://dl.acm.org/doi/10.1145/3510457.3513049}{When Cyber-Physical Systems Meet AI: A Benchmark, an Evaluation, and a Way Forward.}" 44th International Conference on Software Engineering: Software Engineering in Practice. (ICSE 2022, the flagship conference in software engineering, \highlight{Co-first author}, \highlight{CORE Rank A*})
    
    \textbf{Summary}: Presents \strongit{nine pioneering industrial benchmarks} of AI-CPSs across seven domains, \strongit{filling the long-standing gap of systematic evaluation}. Enables multi-view comparisons between AI and classic controllers, showing hybrid controllers superior and motivating AI-aware testing. This work \strongit{laid the benchmarking foundation} for subsequent trustworthy AI-CPS research.
    \item \textbf{[AI-CPS-P2]} Zhenya Zhang, \textbf{Deyun Lyu}, Paolo Arcaini, Lei Ma, Ichiro Hasuo and Jianjun Zhao. ``\href{https://ieeexplore.ieee.org/document/9844247}{FalsifAI: Falsification of AI-Enabled Hybrid Control Systems Guided by Time-Aware Coverage Criteria.}" IEEE Transactions on Software Engineering. (TSE 2023, the flagship journal in software engineering, \highlight{CORE Rank A*}, Impact Factor: 6.5)

    \textbf{Summary}: Proposes \strongit{eight time-aware coverage criteria} tailored for AI-CPS and FalsifAI, the \strongit{first} coverage-guided falsification framework balancing \strongit{exploration} and \strongit{exploitation}. Unlike exploitation-only falsification approaches, it systematically uncovers diverse temporal behaviours. It significantly improves falsification effectiveness and \strongit{pioneered coverage-guided testing for AI-CPSs}, inspiring later testing research.
    \item \textbf{[AI-CPS-P3]} \textbf{Deyun Lyu}, Zhenya Zhang, Paolo Arcaini, Xiao-Yi Zhang, Fuyuki Ishikawa and Jianjun Zhao. ``\href{https://dl.acm.org/doi/10.1145/3705307}{SpectAcle: Fault Localisation of AI-Enabled CPS by Exploiting Sequences of DNN Controller Inferences.}" ACM Transactions on Software Engineering and Methodology. (TOSEM 2025, a top-tier journal in software engineering, \highlight{CORE Rank A*}, Impact Factor: 6.6, \highlight{IPSJ SIG-SE Excellent Research Award})
    
    \textbf{Summary}: Introduces the \strongit{first} fault localisation approach for AI-CPSs. Addressing the lack of per-inference ground truth in sequential control decisions, it exploits temporal sequences of DNN inferences with suspiciousness metrics to \strongit{identify faulty weights}. Evaluated on \strongit{6,067} faulty benchmarks, it effectively localises faults, \strongit{enabling subsequent analysis and repair of AI-CPSs}.
    
    \item \textbf{[AI-CPS-P4]} \textbf{Deyun Lyu}, Yi Li, Zhenya Zhang, Paolo Arcaini, Xiao-Yi Zhang, Fuyuki Ishikawa and Jianjun Zhao. ``\href{https://www.sciencedirect.com/science/article/abs/pii/S0164121225001438}{Fault Localisation of AI-Enabled Cyber-Physical Systems by Exploiting Temporal Neuron Activation.}" Journal of Systems and Software. (JSS 2025, \highlight{CORE Rank A}, Impact Factor: 4.1) 

    \textbf{Summary}: Proposes a neuron-level fault localisation approach for AI-CPSs. By linking temporal neuron activations with system-level specifications, it overcomes the absence of ground truth and identifies faulty neurons. It complements weight-level approaches, \strongit{offering an additional scale of analysis within fault localisation for AI-CPSs}.
    
    \item \textbf{[AI-CPS-P5]} \textbf{Deyun Lyu}, Zhenya Zhang, Paolo Arcaini, Fuyuki Ishikawa, Thomas Laurent and Jianjun Zhao. ``\href{https://dl.acm.org/doi/10.1145/3638529.3654078}{Search-Based Repair of DNN Controllers of AI-Enabled CPS Guided by System-Level Specifications.}"  Proceedings of the Genetic and Evolutionary Computation Conference. (GECCO 2024, a top-tier conference in computational intelligence, \highlight{CORE Rank A})

    \textbf{Summary}: Introduces the \strongit{first} search-based repair approach for DNN controllers of AI-CPSs guided by system-level specifications. Addressing the lack of ground truth, it searches for alternative values of suspicious weights identified through fault localisation, and demonstrates effective, efficient repair that \strongit{strengthens the reliability of AI-CPSs}.
\end{itemize}

\end{document}
